\documentclass[11pt]{article}
\usepackage[a4paper,margin=25mm]{geometry}
\usepackage[T1]{fontenc}
\usepackage{amsmath,amssymb,amsthm,mathtools,microtype}
\usepackage[hidelinks]{hyperref}

\newtheorem{theorem}{Theorem}[section]
\newtheorem{lemma}[theorem]{Lemma}
\newtheorem{corollary}[theorem]{Corollary}
\newtheorem{conjecture}[theorem]{Conjecture}
\theoremstyle{remark}
\newtheorem{remark}[theorem]{Remark}
\DeclareMathOperator{\rank}{rank}
\DeclareMathOperator{\diag}{diag}
\DeclareMathOperator{\tr}{tr}
\newcommand{\R}{\mathbb R}
\newcommand{\CP}{\mathcal{CP}}
\newcommand{\M}{\mathcal M}
\newcommand{\F}{F_2}
\newcommand{\ip}[2]{\langle #1,#2\rangle}
\newcommand{\fro}[1]{\lVert #1\rVert_F}

\title{The Bollob\'as--Nikiforov inequality for nonnegative edge weights}
\author{Gabriel Coutinho, Yinchen Liu, Thom\'as Jung Spier, Quanyu Tang, Shengtong Zhang}
\date{}

\begin{document}
\maketitle

\section{Introduction}

Throughout, graphs are finite and simple, with at least one vertex.
For a graph $G$ on $n$ vertices, let $A_G$ be its adjacency matrix,
$m=|E(G)|$ its number of edges, and $\omega(G)$ its clique number.
We order the adjacency eigenvalues as
$\lambda_1(G)\ge\cdots\ge\lambda_n(G)$.
For matrices of the same size, write
$\ip{B}{C}=\tr(B^{\mathsf T}C)$ and
$\fro{B}^2=\ip{B}{B}$.
The notation $B\succeq0$ means that $B$ is real symmetric and positive
semidefinite, whereas $B\ge0$ means entrywise nonnegativity.
We write $B\circ C$ for the entrywise product and $B^{\circ2}=B\circ B$.
For $t\in\R$, put $t_+=\max\{t,0\}$.
The positive part $B_+$ is taken entrywise, so $(B_+)_{ij}=(B_{ij})_+$.
The all-ones vector and matrix are denoted by $\mathbf1$ and $J$, respectively.
Identity matrices and standard coordinate vectors are denoted by $I$ and $e_i$;
their dimensions will be clear from the context.

The relation between clique number and quadratic optimization goes back to
Motzkin and Straus~\cite{MS}. Their theorem gives
\begin{equation}\label{eq:MS}
 y^{\mathsf T}A_Gy\le
 \left(1-\frac1{\omega(G)}\right)(\mathbf1^{\mathsf T}y)^2
 \qquad(y\ge0).
\end{equation}
This formulation provided a new proof of Tur\'an's theorem and a natural starting
point for spectral extensions. Wilf~\cite{Wilf} obtained spectral lower bounds
for the clique number, including
$\lambda_1(G)\le (1-1/\omega(G))n$.
Nikiforov~\cite{Nikiforov} subsequently proved the edge-count bound
\[
 \lambda_1(G)^2\le 2\left(1-\frac1{\omega(G)}\right)m.
\]
Bollob\'as and Nikiforov~\cite{BN} proposed that the same right-hand side controls
the squares of the two largest eigenvalues, provided the graph is not complete.

\begin{conjecture}\label{conj:BN}
Every noncomplete graph $G$ on at least two vertices satisfies
\[
 \lambda_1(G)^2+\lambda_2(G)^2
 \le 2\left(1-\frac1{\omega(G)}\right)|E(G)|.
\]
\end{conjecture}

Several important cases have been established. The theorem of Ando and
Lin~\cite{AL} on the sum of squares of positive eigenvalues implies the
conjecture when the chromatic number equals the clique number.
Lin, Ning, and Wu~\cite{LNW} proved it for triangle-free graphs and determined
the equality cases in that setting. Zhang~\cite{Zhang} proved the regular case.
Kumar and Pragada~\cite{KP} obtained results for graphs with few triangles,
while Zeng and Zhang~\cite{ZZ} treated line graphs and further classes with few
triangles. Liu and Bu~\cite{LB} established an asymptotically almost sure
version for random graphs. The survey of Liu and Ning~\cite{LN} places the
conjecture among related questions in spectral graph theory.

Coutinho, Spier, and Zhang~\cite{CSZ} developed a conic optimization approach
to sums of squares of eigenvalues and obtained versions of
Conjecture~\ref{conj:BN} with weaker constants. Their work emphasizes the
difference between complete positivity and the conjunction of positive
semidefiniteness and entrywise nonnegativity. It also introduces a rank-two
optimization problem whose conjectured value is the clique number.
We use this viewpoint, together with total nonnegativity and Schur complements,
to prove Conjecture~\ref{conj:BN}.

For a real symmetric matrix $B$, let $\F(B)$ be the sum of the squares of its
two largest positive eigenvalues, with missing terms replaced by zero.
Our spectral result allows arbitrary nonnegative edge weights.

\begin{theorem}\label{thm:weighted}
Let $G$ be a graph, and let $B$ be a symmetric entrywise nonnegative matrix
with zero diagonal such that $B_{ij}=0$ whenever $\{i,j\}\notin E(G)$.
Then
\begin{equation}\label{eq:weighted}
 \F(B)\le\left(1-\frac1{\omega(G)}\right)\fro{B}^2.
\end{equation}
\end{theorem}

For each fixed graph with an edge, the coefficient in~\eqref{eq:weighted}
is best possible: one may take the adjacency matrix of a maximum clique and
extend it by zeros. A noncomplete adjacency matrix has a nonnegative second
largest eigenvalue, so Theorem~\ref{thm:weighted} gives
Conjecture~\ref{conj:BN}. We supply the details in Section~\ref{sec:spectral}.

A symmetric matrix $C$ is \emph{completely positive} if
$C=\sum_{a=1}^q p_ap_a^{\mathsf T}$ for finitely many entrywise nonnegative
vectors $p_a$. We denote this cone by $\CP_n$ in dimension $n$.
For a positive semidefinite matrix $X$, define
\begin{equation}\label{eq:matrix}
 \M(X)=X^{\circ2}+
 \sum_{i<j:X_{ij}<0}X_{ij}^2(e_i-e_j)(e_i-e_j)^{\mathsf T}.
\end{equation}
This matrix is positive semidefinite and entrywise nonnegative.
The following stronger property is the main matrix result.

\begin{theorem}\label{thm:matrix}
Suppose $X_{ij}=z_i^{\mathsf T}z_j$, where $z_1,\ldots,z_n\in\R^2$.
If there is a nonzero vector $w\in\R^2$ such that
$w^{\mathsf T}z_i\ge0$ for every $i$, then $\M(X)$ is completely positive.
\end{theorem}

Thus a common closed half-plane, with boundary through the origin, suffices.
In the spectral application, this condition is supplied by a nonnegative
Perron eigenvector. In particular, irreducibility is not needed.
After obtaining Theorem~\ref{thm:weighted}, a variational argument removes
the half-plane restriction from the resulting scalar inequality.

\begin{theorem}\label{thm:gram}
For every graph $G$ and every real positive semidefinite matrix $X$ of rank
at most two,
\begin{equation}\label{eq:gram}
 \sum_{i,j=1}^n(A_G)_{ij}(X_{ij})_+^2
 \le\left(1-\frac1{\omega(G)}\right)\fro{X}^2.
\end{equation}
\end{theorem}

The sum in~\eqref{eq:gram} counts each edge twice.
This result also settles the two equivalent conjectures numbered 2 and 3
in~\cite{CSZ}, without their restriction to sufficiently large clique number;
see Section~\ref{subsec:conic}. The rank restriction remains essential to our
argument. We do not address the corresponding bound involving up to
$\omega(G)$ positive eigenvalues proposed by Elphick, Linz, and
Wocjan~\cite{ELW}.

The proof of Theorem~\ref{thm:matrix} has two parts. An ordered elimination
produces nonnegative columns. The remaining Schur complement is the sum of
a graph Laplacian and a matrix with an explicit three-column nonnegative
factorization. A standard operation on completely positive matrices then
recovers the original matrix. The sign calculation in the three-column
factorization reduces to total nonnegativity of two rectangular matrices.

Section~\ref{sec:prelim} records the facts about complete positivity and
total nonnegativity that we use. Section~\ref{sec:kernel} proves the
three-column factorization. Section~\ref{sec:matrix} proves
Theorem~\ref{thm:matrix}. Section~\ref{sec:spectral} derives the spectral
inequality, Theorem~\ref{thm:gram}, and the rank-two optimization identity.

\section{Complete positivity and total nonnegativity}\label{sec:prelim}

We first record a standard operation on completely positive matrices;
compare~\cite{SMB} and the proof of Theorem 29 in~\cite{CSZ}.
The short proof specifies exactly how much an off-diagonal entry may decrease.

\begin{lemma}\label{lem:pair}
If $C\in\CP_n$, $i\ne j$, and $0\le h\le C_{ij}$, then
\[
 C+h(e_i-e_j)(e_i-e_j)^{\mathsf T}\in\CP_n.
\]
\end{lemma}
\begin{proof}
First let $C=pp^{\mathsf T}$ with $p\ge0$, and write $a=p_i$, $b=p_j$.
If $ab>0$, form $u$ and $v$ by replacing $(p_i,p_j)$ with $(a+b,0)$ and
$(0,a+b)$, respectively, and leaving the other coordinates unchanged.
Then
\[
 pp^{\mathsf T}+ab(e_i-e_j)(e_i-e_j)^{\mathsf T}
 =\frac{a}{a+b}uu^{\mathsf T}+\frac{b}{a+b}vv^{\mathsf T}.
\]
Convex interpolation proves the assertion for $0\le h\le ab$.
When $ab=0$, only $h=0$ is allowed.
For a decomposition $C=\sum_a p_ap_a^{\mathsf T}$, distribute $h$ among
numbers $h_a$ satisfying $0\le h_a\le(p_a)_i(p_a)_j$ and $\sum_a h_a=h$.
Such a distribution exists because the upper bounds sum to $C_{ij}$.
Adding the rank-one conclusions finishes the proof.
\end{proof}

A \emph{weighted graph Laplacian} means a matrix of the form
$L=\sum_{i<j}\ell_{ij}(e_i-e_j)(e_i-e_j)^{\mathsf T}$ with $\ell_{ij}\ge0$.

\begin{corollary}\label{cor:laplacian}
If $C$ is completely positive, $L$ is a weighted graph Laplacian, and
$C+L\ge0$ entrywise, then $C+L$ is completely positive.
\end{corollary}
\begin{proof}
For each pair, entrywise nonnegativity gives $\ell_{ij}\le C_{ij}$.
Apply Lemma~\ref{lem:pair} successively. An update on another pair does not
change the $(i,j)$ entry, so every application satisfies its hypothesis.
\end{proof}

We also use closedness of $\CP_n$ and its closure under taking principal
submatrices. The latter is immediate. For the former, conic
Carath\'eodory gives a factorization with at most $n(n+1)/2$ nonnegative
columns. If a sequence of completely positive matrices converges, its traces
are bounded, and hence so are these columns. Padding with zero columns and
passing to a convergent subsequence proves closedness.
These are standard properties of the cone; see~\cite{SMB}.

A matrix is \emph{totally nonnegative} if every square minor is nonnegative.
For a kernel on ordered sets, the definition requires nonnegative determinants
on increasing lists of row and column arguments. We use the following two
elementary instances of this notion; general background is given in~\cite{FJ}.

\begin{lemma}\label{lem:truncated}
The kernel $K(a,t)=(t-a)_+^2$ is totally nonnegative.
\end{lemma}
\begin{proof}
The step kernel $H(a,t)=\mathbf1_{\{a\le t\}}$ has ordered minors equal to
zero or one. Indeed, its rows have nested terminal intervals of ones;
subtracting each successive row from the preceding row verifies the claim.
The continuous Cauchy--Binet identity implies that composition preserves total
nonnegativity: an ordered minor of a composed kernel is the integral over
increasing intermediate arguments of a product of ordered minors.
All integrations here may be restricted to a bounded interval containing the
external arguments. Since
\[
 \int\!\!\int H(a,s)H(s,u)H(u,t)\,ds\,du=\tfrac12(t-a)_+^2,
\]
two compositions and positive scaling give the result.
\end{proof}

\begin{lemma}\label{lem:convex}
Let $f:[0,\infty)\to[0,\infty)$ be convex with $f(0)=0$.
For $0<t_1\le\cdots\le t_k$, the matrix whose $i$th row is
$(1,t_i,f(t_i))$ is totally nonnegative.
\end{lemma}
\begin{proof}
Convexity and $f(0)=0$ imply that $f(t)/t$ is nondecreasing for $t>0$.
Since $f\ge0$, the function $f$ is also nondecreasing. These observations
verify all minors of orders one and two.
For $t<u<v$, the minor of order three is
\[
 (u-t)\bigl(f(v)-f(t)\bigr)
 -(v-t)\bigl(f(u)-f(t)\bigr),
\]
which is nonnegative by convexity. Repeated arguments give zero minors.
There are only three columns, so no larger minors need to be considered.
\end{proof}

\section{A three-column nonnegative factorization}\label{sec:kernel}

Fix $0<t_1\le\cdots\le t_k$, $q_i>0$, and $\gamma>0$.
In this section the coordinates of $\R^3$ are numbered $0,1,2$.
For $x\ge0$, set
\[
\begin{gathered}
 v(t)=(1,-\sqrt2t,t^2)^{\mathsf T},\qquad
 b(x)=(x^2,\sqrt2x,1)^{\mathsf T},\qquad a_i(x)=(x-t_i)_+^2,\\
 \mathcal A=I_3+\sum_iq_iv(t_i)v(t_i)^{\mathsf T},\qquad V=\mathcal A e_0,\\
 h(x)=\sum_iq_i a_i(x),\qquad h_1(x)=\sum_iq_it_i a_i(x),\\
 \widehat b(x)=b(x)+\sum_iq_i a_i(x)v(t_i),\qquad
 U(x)=\widehat b(x)+\frac{h(x)}\gamma V,\\
 \mathcal K=\left(\mathcal A+\frac{VV^{\mathsf T}}\gamma\right)^{-1}.
\end{gathered}
\]
Write $m_j=\sum_iq_it_i^j$ for $0\le j\le4$, and define
\[
 a_0=1+m_0,\qquad D_2=a_0(1+2m_2)-2m_1^2,\qquad
 \Delta=\det\mathcal A.
\]
Both $D_2$ and $\Delta$ are positive, since they are principal minors of the
positive definite matrix $\mathcal A$.
Finally, put
\begin{equation}\label{eq:functions}
\begin{aligned}
 N(x)&=\Delta e_2^{\mathsf T}\mathcal A^{-1}\widehat b(x),\\
 P(x)&=a_0x+m_1x^2+m_1h(x)-a_0h_1(x),\\
 Z(x)&=\gamma x^2+(\gamma+a_0)h(x).
\end{aligned}
\end{equation}

\begin{lemma}\label{lem:kernel}
The functions in~\eqref{eq:functions} satisfy $N(x)\ge1$, $P(x)\ge0$, and
$Z(x)\ge0$ for $x\ge0$, and
\begin{equation}\label{eq:factor}
 U(x)^{\mathsf T}\mathcal K U(y)
 =\frac{N(x)N(y)}{\Delta D_2}
  +\frac{2P(x)P(y)}{a_0D_2}
  +\frac{Z(x)Z(y)}{\gamma a_0(\gamma+a_0)}
 \qquad(x,y\ge0).
\end{equation}
\end{lemma}
\begin{proof}
Since $\mathcal A^{-1}V=e_0$, the Sherman--Morrison identity gives
\[
 \mathcal K=\mathcal A^{-1}-\frac{e_0e_0^{\mathsf T}}{\gamma+a_0}.
\]
For clarity, completing squares in the order $2,1,0$ yields, for arbitrary
$z,w\in\R^3$,
\begin{equation}\label{eq:bilinear}
\begin{aligned}
 z^{\mathsf T}\mathcal K w
 ={}&\frac{(\Delta e_2^{\mathsf T}\mathcal A^{-1}z)
              (\Delta e_2^{\mathsf T}\mathcal A^{-1}w)}{\Delta D_2}\\
 &+\frac{(a_0z_1+\sqrt2m_1z_0)(a_0w_1+\sqrt2m_1w_0)}{a_0D_2}
  +\frac{\gamma z_0w_0}{a_0(\gamma+a_0)}.
\end{aligned}
\end{equation}
Here the first pivot is $\mathcal K_{22}=D_2/\Delta$.
After eliminating coordinate $2$, the remaining block is
$((\mathcal K^{-1})_{\{0,1\},\{0,1\}})^{-1}$; its $(1,1)$ entry is
$a_0/D_2$, and the ratio of its $(1,0)$ and $(1,1)$ entries is
$\sqrt2m_1/a_0$. The last pivot is
$1/(\mathcal K^{-1})_{00}=\gamma/[a_0(\gamma+a_0)]$.
These identities verify~\eqref{eq:bilinear}.
Substitute $z=U(x)$ and $w=U(y)$, and use
\[
 \Delta e_2^{\mathsf T}\mathcal A^{-1}U(x)=N(x),\qquad
 a_0U_1(x)+\sqrt2m_1U_0(x)=\sqrt2P(x),\qquad U_0(x)=Z(x)/\gamma.
\]
This proves~\eqref{eq:factor}.

Clearly $Z(x)\ge0$. Expanding $P$ gives
\[
\begin{aligned}
 P(x)={}&a_0x+\sum_iq_it_i\bigl[x^2-(x-t_i)_+^2\bigr]\\
 &+\sum_{i<j}q_iq_j(t_j-t_i)
       \bigl[(x-t_i)_+^2-(x-t_j)_+^2\bigr].
\end{aligned}
\]
Each summand is nonnegative, so $P(x)\ge0$.

It remains to prove $N(x)\ge1$.
Let $D(x)$ be the matrix with columns $e_0,e_1,b(x)$, and put
$c_i(x)=(1,-\sqrt2t_i,a_i(x))^{\mathsf T}$.
Cramer's rule gives
\[
 N(x)=\det[\mathcal A e_0,\mathcal A e_1,\widehat b(x)]
     =\det\left(D(x)+\sum_iq_iv(t_i)c_i(x)^{\mathsf T}\right).
\]
Since $\det D(x)=1$, Sylvester's determinant identity implies
\begin{equation}\label{eq:Ndet}
 N(x)=\det(I_k+QW(x)),\qquad
 Q=\diag(q_1,\ldots,q_k),\qquad
 W(x)_{ij}=c_i(x)^{\mathsf T}D(x)^{-1}v(t_j).
\end{equation}
Direct multiplication gives
\begin{equation}\label{eq:W}
 W(x)_{ij}=1+2t_it_j+f_x(t_i)t_j^2,\qquad
 f_x(t)=t^2-(t-x)_+^2.
\end{equation}
On $[0,\infty)$, the function $f_x$ is nonnegative and convex, with
$f_x(0)=0$: explicitly, it equals $t^2$ for $t\le x$ and $2xt-x^2$ for
$t\ge x$, and its derivative is $2\min\{t,x\}$.
By Lemma~\ref{lem:convex}, the matrices with rows
$(1,t_i,f_x(t_i))$ and $(1,2t_i,t_i^2)$ are totally nonnegative.
Equation~\eqref{eq:W} expresses $W(x)$ as the product of the first matrix
and the transpose of the second. The Cauchy--Binet formula therefore shows
that every minor of $W(x)$ is nonnegative.
Expanding~\eqref{eq:Ndet} in principal minors now gives
\[
 N(x)=\sum_{S\subseteq\{1,\ldots,k\}}
       \left(\prod_{i\in S}q_i\right)\det W(x)[S,S]\ge1,
\]
where the empty minor is defined to be $1$.
This completes the proof.
\end{proof}

In particular, for any finite set of nonnegative arguments, the matrix with
entries $U(x)^{\mathsf T}\mathcal K U(y)$ is completely positive with at most
three nonnegative columns. We will use this consequence of
Lemma~\ref{lem:kernel} below.

\section{A completely positive matrix from a planar Gram representation}
\label{sec:matrix}

The map in~\eqref{eq:matrix} is continuous, and direct inspection gives
\begin{equation}\label{eq:mass}
 \M(X)_{ij}=(X_{ij})_+^2\quad(i\ne j),\qquad
 \M(X)\mathbf1=X^{\circ2}\mathbf1,\qquad
 \ip{J}{\M(X)}=\fro{X}^2.
\end{equation}
We first work with a normalized configuration and later pass to limits.
Consider
\begin{equation}\label{eq:coordinates}
 z_0=(1,0),\qquad z_i=\sqrt{s_i}(-1,t_i)\ (1\le i\le k),\qquad
 y_j=\sqrt{\rho_j}(x_j,1)\ (1\le j\le p),
\end{equation}
where $k\ge1$, $s_i,\rho_j>0$, $0<t_1\le\cdots\le t_k$, and $x_j\ge0$.
Let $X$ be their Gram matrix and write $M=\M(X)$.
Use $E=\{0,1,\ldots,k\}$ for the $z$-indices and
$T=\{\bar1,\ldots,\bar p\}$ for the $y$-indices.
Set
\[
 a_i(x)=(x-t_i)_+^2,\qquad H_i=\sum_j\rho_j a_i(x_j),\qquad
 d_i=1+H_i,\qquad \sigma=\sum_i s_i.
\]
The entries needed for elimination are
\begin{equation}\label{eq:blocks}
\begin{aligned}
 M_{00}&=1+\sigma,& M_{0i}&=0,& M_{0\bar j}&=\rho_jx_j^2,\\
 M_{ih}&=s_is_h(1+t_it_h)^2+\delta_{ih}s_id_i,
 &&1\le i,h\le k,\\
 M_{i\bar j}&=s_i\rho_j(t_i-x_j)_+^2,
 &&1\le i\le k,\ 1\le j\le p.
\end{aligned}
\end{equation}
Here $\delta_{ih}$ is $1$ when $i=h$ and $0$ otherwise.

\begin{lemma}\label{lem:elimination}
For the configuration~\eqref{eq:coordinates}, there is a completely positive
matrix $C_0$ such that
\begin{equation}\label{eq:extraction}
 M=C_0+\widetilde R,\qquad
 R=M_{TT}-M_{TE}M_{EE}^{-1}M_{ET}\succeq0,
\end{equation}
where a tilde denotes extension by zero outside the rows and columns in $T$.
\end{lemma}
\begin{proof}
We eliminate the indices $0,1,\ldots,k$ in that order by symmetric Gaussian
elimination. The pivot at $0$ is positive and its column is nonnegative.
By~\eqref{eq:blocks}, this step changes neither the block on $\{1,\ldots,k\}$
nor its coupling to $T$.
The former block is the sum of a Gram matrix and the positive diagonal matrix
$\diag(s_1d_1,\ldots,s_kd_k)$, so all subsequent pivots are positive.

It suffices to show that each of the remaining elimination columns is
nonnegative. At pivot $j$, let $I=(1,\ldots,j-1)$ and let $\alpha$ be a
later index. The residual entry in row $\alpha$ and column $j$ is
\[
 \frac{\det M[I,\alpha;I,j]}{\det M[I,I]}.
\]
The notation in the numerator specifies the ordered row list $(I,\alpha)$
and column list $(I,j)$; for $I$ empty the denominator is $1$.
These entries involve only the unchanged block and its couplings.
Expand the numerator in the positive diagonal terms indexed by $I$.
Deleting the same row and column positions contributes positive cofactor sign.
After removing the positive row and column factors $s_i$ and $\rho_j$,
every resulting term is a minor whose ordinary rows have entries
$(1+t_at_b)^2$. There may also be a last row with entries $(t_b-x)_+^2$,
where $x\ge0$.

For $s,t>0$,
\[
 (1+st)^2=s^2\bigl(t-(-1/s)\bigr)_+^2.
\]
Thus these are, up to positive row factors, ordered minors of the kernel in
Lemma~\ref{lem:truncated}. The row arguments $-1/t_a$ increase with $t_a$.
A later $z$-index preserves this order, and a last argument $x\ge0$ follows
all the negative row arguments. The column arguments are also increasing.
All the terms are therefore nonnegative.

Dividing each elimination column by the square root of its positive pivot
produces a nonnegative vector. The sum of their outer products is a completely
positive matrix $C_0$. The remaining block is the Schur complement in
\eqref{eq:extraction}, which is positive semidefinite because $M\succeq0$.
\end{proof}

We now identify the remaining block. Put
\begin{equation}\label{eq:qgamma}
 q_i=\frac{s_i}{d_i},\qquad
 \gamma=\sigma-\sum_iq_i=\sum_i\frac{s_iH_i}{1+H_i}.
\end{equation}
For the moment assume $\gamma>0$.
Use the notation of Section~\ref{sec:kernel} with these $q_i$ and $\gamma$.
Let $F$ have rows $e_0^{\mathsf T}$ at index $0$,
$s_iv(t_i)^{\mathsf T}$ at index $i$, and
$\rho_jb(x_j)^{\mathsf T}$ at index $\bar j$.
These are the degree-two Gram features, so
\begin{equation}\label{eq:features}
 M=FF^{\mathsf T}+L,
\end{equation}
where $L$ is the Laplacian of the pairs with negative original inner product.
Its only possibly nonzero edge weights are $s_i$ on $0i$ and
$s_i\rho_j a_i(x_j)$ on $i\bar j$.
Writing $s=(s_1,\ldots,s_k)^{\mathsf T}$, we have
\[
 L_{EE}=\begin{pmatrix}
 \sigma&-s^{\mathsf T}\\
 -s&\diag(s_1d_1,\ldots,s_kd_k)
 \end{pmatrix}\succ0.
\]
The Schur complement of the lower-right block is $\gamma$, and
\begin{equation}\label{eq:Linverse}
 L_{EE}^{-1}=\diag\left(0,\frac1{s_1d_1},\ldots,\frac1{s_kd_k}\right)
             +\frac1\gamma ww^{\mathsf T},\qquad
 w=(1,d_1^{-1},\ldots,d_k^{-1})^{\mathsf T}.
\end{equation}

Define
\[
 L_{\rm red}=L_{TT}-L_{TE}L_{EE}^{-1}L_{ET},\qquad
 \mathcal U=F_T-L_{TE}L_{EE}^{-1}F_E.
\]
A block elimination in~\eqref{eq:features}, followed by the Woodbury identity,
gives
\begin{equation}\label{eq:Schur}
 R=L_{\rm red}+\mathcal U
       (I_3+F_E^{\mathsf T}L_{EE}^{-1}F_E)^{-1}\mathcal U^{\mathsf T}.
\end{equation}
Indeed, the block elimination that replaces $L$ by
$\diag(L_{EE},L_{\rm red})$ replaces the feature rows by $(F_E,\mathcal U)$.
The Schur complement of the first block is then
\[
 L_{\rm red}+\mathcal U
 \bigl[I_3-F_E^{\mathsf T}(L_{EE}+F_EF_E^{\mathsf T})^{-1}F_E\bigr]
 \mathcal U^{\mathsf T},
\]
and the bracket equals the inverse appearing in~\eqref{eq:Schur}.

By~\eqref{eq:Linverse}, the center row of $L_{EE}^{-1}F_E$ is
$V^{\mathsf T}/\gamma$, and its $i$th row is
$(v(t_i)+V/\gamma)^{\mathsf T}/d_i$. Hence
\[
 I_3+F_E^{\mathsf T}L_{EE}^{-1}F_E
 =\mathcal A+VV^{\mathsf T}/\gamma,\qquad
 \mathcal U_{\bar j}=\rho_jU(x_j)^{\mathsf T}.
\]
The second term on the right of~\eqref{eq:Schur} therefore has entries
$\rho_j\rho_\ell U(x_j)^{\mathsf T}\mathcal K U(x_\ell)$.
It is completely positive by Lemma~\ref{lem:kernel}; denote it by $C_1$.

The matrix $L_{\rm red}$ is a weighted graph Laplacian. This is the usual
Schur-complement property called Kron reduction; see~\cite{DB}.
In this instance it also follows directly from
\[
 -(L_{\rm red})_{\bar j\bar\ell}
 =\rho_j\rho_\ell\left[
   \sum_iq_i a_i(x_j)a_i(x_\ell)+\frac{h(x_j)h(x_\ell)}\gamma
   \right]\ge0\qquad(j\ne\ell)
\]
and $L_{\rm red}\mathbf1=0$.
The latter identity follows from $L\mathbf1=0$ by solving its $E$-rows for
the vector on $E$ and substituting into its $T$-rows.
Combining Lemma~\ref{lem:elimination} with~\eqref{eq:Schur}, we obtain
\[
 M=C_0+\widetilde C_1+\widetilde L_{\rm red}.
\]
The first two terms are completely positive, and $M\ge0$ entrywise.
Corollary~\ref{cor:laplacian} proves that $M$ is completely positive.
Notice that we do not need the Schur complement $R$ itself to be entrywise
nonnegative.

The case $\gamma=0$ in~\eqref{eq:qgamma} follows by adjoining one small vector.
Choose $x_*>\max_i t_i$ and add
$y_*=\sqrt\varepsilon(x_*,1)$, where $\varepsilon>0$.
Every $H_i$ in the enlarged configuration is then positive, so the preceding
argument applies. The principal submatrix of its completely positive matrix
on the original indices converges to $M$ as $\varepsilon\downarrow0$.
Indeed, the new vector changes the old entries only by diagonal contributions
of order $\varepsilon$ in~\eqref{eq:matrix}.
Closedness of the completely positive cone proves the assertion for the
original configuration. This argument also covers $p=0$.

We can now prove Theorem~\ref{thm:matrix}.
\begin{proof}
First suppose that all nonzero Gram vectors lie in a common open half-plane,
and remove zero vectors. If no pair has negative inner product, their angular
span is at most $\pi/2$. After a rotation they have nonnegative coordinates,
so $X=pp^{\mathsf T}+qq^{\mathsf T}$ with $p,q\ge0$. Therefore
\[
 \M(X)=X^{\circ2}
 =p^{\circ2}(p^{\circ2})^{\mathsf T}
  +q^{\circ2}(q^{\circ2})^{\mathsf T}
  +2(p\circ q)(p\circ q)^{\mathsf T}
\]
is completely positive.

Otherwise, first assume the minimum angular direction is unique.
Rotate and multiply all vectors by a common positive scalar to make that
vector $(1,0)$. All other angles lie in $(0,\pi)$.
Separating the vectors by the sign of their first coordinate gives
\eqref{eq:coordinates}, with $k\ge1$, after ordering the parameters $t_i$.
The result just proved applies. The common rescaling causes no difficulty,
since replacing $X$ by $c^2X$ replaces $\M(X)$ by $c^4\M(X)$.
Tied minimum directions are handled by a perturbation within the open
half-plane and closedness of $\CP_n$.

Finally, let $w$ be a unit vector defining the closed half-plane in the
statement. Replacing each $z_i$ by $z_i+\varepsilon w$ puts every vector in
the corresponding open half-plane. Their Gram matrices converge to $X$.
Continuity of~\eqref{eq:matrix} and closedness of $\CP_n$ give the result.
\end{proof}

\section{Spectral inequalities and rank-two optimization}\label{sec:spectral}

The completely positive form of Motzkin--Straus is
\begin{equation}\label{eq:CPMS}
 \ip{A_G}{C}\le\left(1-\frac1{\omega(G)}\right)\ip{J}{C}
 \qquad(C\in\CP_n).
\end{equation}
It follows by applying~\eqref{eq:MS} to every column in a nonnegative
factorization of $C$ and adding. This formulation is classical in copositive
optimization; see de Klerk and Pasechnik~\cite{dKP} and~\cite{CSZ}.

We first prove Theorem~\ref{thm:weighted}.
\begin{proof}
The case $B=0$ is immediate. Otherwise set $r=\omega(G)\ge2$, and let
$a>0$ be the largest eigenvalue of $B$. Take a unit nonnegative Perron
eigenvector $u$. If $B$ has a second positive eigenvalue $b$, choose a unit
eigenvector $v$ orthogonal to $u$ and put
$X=auu^{\mathsf T}+bvv^{\mathsf T}$.
If there is no second positive eigenvalue, put $X=auu^{\mathsf T}$.
In either case,
\[
 \ip{B}{X}=\fro{X}^2=\F(B)=:S>0.
\]
The planar Gram vectors of $X$ can be taken to be
$(\sqrt a\,u_i,\sqrt b\,v_i)$, with second coordinate zero when the second
term is absent. They lie in the closed right half-plane, even when $B$ is
reducible. By Theorem~\ref{thm:matrix}, $\M(X)$ is completely positive.
Equations~\eqref{eq:mass} and~\eqref{eq:CPMS} give
\[
 \sum_{i,j}(A_G)_{ij}(X_{ij})_+^2\le(1-1/r)S.
\]
Using the support and nonnegativity of $B$, followed by Cauchy--Schwarz,
we obtain
\[
 S=\ip{B}{X}\le\ip{B}{A_G\circ X_+}
 \le\fro{B}\left(\sum_{i,j}(A_G)_{ij}(X_{ij})_+^2\right)^{1/2}
 \le\fro{B}\sqrt{(1-1/r)S}.
\]
Squaring and dividing by $S$ proves the theorem.
\end{proof}

For sharpness in~\eqref{eq:weighted}, take the adjacency matrix of a maximum
clique of size $r\ge2$ and extend it by zeros. Its only positive eigenvalue is
$r-1$, and its squared Frobenius norm is $r(r-1)$.

To obtain Conjecture~\ref{conj:BN}, observe that a nonedge gives a zero
$2\times2$ principal submatrix of $A_G$.
Eigenvalue interlacing yields $\lambda_2(G)\ge0$.
Thus $\F(A_G)=\lambda_1(G)^2+\lambda_2(G)^2$, while
$\fro{A_G}^2=2|E(G)|$. Theorem~\ref{thm:weighted} proves the conjecture;
edgeless graphs are included.

The following elementary variational inequality will allow arbitrary planar
Gram vectors.

\begin{lemma}\label{lem:variational}
If $B$ is real symmetric and $X\succeq0$ has rank at most two, then
\[
 \ip{B}{X}\le\sqrt{\F(B)}\,\fro{X}.
\]
\end{lemma}
\begin{proof}
The case of matrices of order one follows directly, so assume the order is
at least two. Write
$X=\alpha uu^{\mathsf T}+\beta vv^{\mathsf T}$, where $u,v$ are orthonormal
and $\alpha\ge\beta\ge0$, allowing a zero coefficient.
If $\mu_1\ge\mu_2$ are the two largest eigenvalues of $B$, the Rayleigh
principle and the two-dimensional variational principle give
\[
 \ip{B}{X}
 =(\alpha-\beta)u^{\mathsf T}Bu+
   \beta(u^{\mathsf T}Bu+v^{\mathsf T}Bv)
 \le\alpha\mu_1+\beta\mu_2.
\]
For completeness, the second variational principle follows by expanding the
orthogonal projection $uu^{\mathsf T}+vv^{\mathsf T}$ in an eigenbasis of $B$:
its diagonal entries lie in $[0,1]$ and sum to $2$.
Replacing $\mu_1,\mu_2$ by their positive parts and applying Cauchy--Schwarz
proves the assertion.
\end{proof}

We now prove Theorem~\ref{thm:gram}.
\begin{proof}
Set $B=A_G\circ X_+$ and $T=\fro{B}^2$.
Then $B$ satisfies the hypotheses of Theorem~\ref{thm:weighted}, and
$\ip{B}{X}=T$.
If $T=0$, the assertion is immediate. Otherwise, Theorem~\ref{thm:weighted}
and Lemma~\ref{lem:variational} give
\[
 T\le\sqrt{\F(B)}\,\fro{X}
 \le\sqrt{\left(1-\frac1{\omega(G)}\right)T}\,\fro{X}.
\]
Squaring and dividing by $T$ proves~\eqref{eq:gram}.
\end{proof}

\subsection{The rank-two optimization problem}\label{subsec:conic}

Let $\overline G$ be the complement of $G$, so
$A_{\overline G}=J-I-A_G$.
Following~\cite{CSZ}, define
\begin{equation}\label{eq:parameter}
 \chi''_{\mathrm{vec},3}(G)=
 \max\left\{\ip{J}{X^{\circ2}}:
 \begin{array}{l}
 X\succeq0,\quad\rank X\le2,\\
 \ip{I+A_{\overline G}}{X^{\circ2}}=1,\\
 X_{ij}\ge0\quad\text{for }\{i,j\}\in E(G)
 \end{array}\right\}.
\end{equation}
This maximum is well defined. The feasible set is nonempty, and the
normalization implies $\sum_iX_{ii}^2\le1$.
Positive semidefiniteness then bounds every entry by
$|X_{ij}|^2\le X_{ii}X_{jj}\le1$. All constraints are closed, so the
feasible set is compact.

Conjecture 2 of Coutinho, Spier, and Zhang~\cite{CSZ} is the following.
\begin{conjecture}\label{conj:CSZ}
There is an integer $r_0$ such that
$\chi''_{\mathrm{vec},3}(G)=\omega(G)$ whenever $\omega(G)\ge r_0$.
\end{conjecture}

Theorem~\ref{thm:gram} proves a stronger statement.
\begin{corollary}\label{cor:parameter}
Every graph $G$ satisfies $\chi''_{\mathrm{vec},3}(G)=\omega(G)$.
\end{corollary}
\begin{proof}
Let $r=\omega(G)$ and let $X$ be feasible in~\eqref{eq:parameter}.
Since $X$ is nonnegative on edges, Theorem~\ref{thm:gram} gives
\[
 \ip{A_G}{X^{\circ2}}\le(1-1/r)\fro{X}^2.
\]
The normalization therefore implies
\[
 1=\fro{X}^2-\ip{A_G}{X^{\circ2}}\ge\frac1r\fro{X}^2,
\]
so the objective is at most $r$.
For the reverse inequality, let $q$ be the indicator vector of a clique of
size $r$, and take $X=qq^{\mathsf T}/\sqrt r$.
This matrix is feasible and has objective value $r$.
\end{proof}

This proves Conjecture~\ref{conj:CSZ} without a lower bound on the clique
number. The equivalent planar-vector formulation, Conjecture 3 in~\cite{CSZ},
follows by writing $X_{ij}=z_i^{\mathsf T}z_j$ in~\eqref{eq:gram}.
More explicitly, for arbitrary $z_1,\ldots,z_n\in\R^2$ we have
\[
 \sum_{\{i,j\}\in E(G)}(z_i^{\mathsf T}z_j)_+^2
 \le\frac12\left(1-\frac1{\omega(G)}\right)
             \sum_{i,j=1}^n(z_i^{\mathsf T}z_j)^2.
\]
When inner products on edges are nonnegative, the positive parts can be
omitted. No common half-plane assumption is needed for this scalar inequality.
Corollary~\ref{cor:parameter} concerns this rank-two problem; it does not assert
complete positivity of~\eqref{eq:matrix} for arbitrary planar configurations.

\begin{thebibliography}{99}

\bibitem{AL}
T.~Ando and M.~Lin,
Proof of a conjectured lower bound on the chromatic number of a graph,
\emph{Linear Algebra Appl.} \textbf{485} (2015), 480--484.

\bibitem{BN}
B.~Bollob\'as and V.~Nikiforov,
Cliques and the spectral radius,
\emph{J. Combin. Theory Ser. B} \textbf{97} (2007), 859--865.

\bibitem{CSZ}
G.~Coutinho, T.~J.~Spier, and S.~Zhang,
Conic programming to understand sums of squares of eigenvalues of graphs,
arXiv:2411.08184 (2024).

\bibitem{dKP}
E.~de Klerk and D.~V.~Pasechnik,
Approximation of the stability number of a graph via copositive programming,
\emph{SIAM J. Optim.} \textbf{12} (2002), 875--892.

\bibitem{DB}
F.~D\"orfler and F.~Bullo,
Kron reduction of graphs with applications to electrical networks,
\emph{IEEE Trans. Circuits Syst. I Regul. Pap.} \textbf{60} (2013), 150--163.

\bibitem{ELW}
C.~Elphick, W.~Linz, and P.~Wocjan,
Two conjectured strengthenings of Tur\'an's theorem,
\emph{Linear Algebra Appl.} \textbf{684} (2024), 23--36.

\bibitem{FJ}
S.~M.~Fallat and C.~R.~Johnson,
\emph{Totally Nonnegative Matrices},
Princeton University Press, Princeton, NJ, 2011.

\bibitem{KP}
H.~Kumar and S.~Pragada,
Bollob\'as--Nikiforov conjecture for graphs with not so many triangles,
\emph{Linear Algebra Appl.} \textbf{727} (2025), 1--9.

\bibitem{LNW}
H.~Lin, B.~Ning, and B.~Wu,
Eigenvalues and triangles in graphs,
\emph{Combin. Probab. Comput.} \textbf{30} (2021), 258--270.

\bibitem{LB}
C.~Liu and C.~Bu,
Bollob\'as--Nikiforov conjecture holds asymptotically almost surely,
arXiv:2501.07137 (2025).

\bibitem{LN}
L.~Liu and B.~Ning,
Unsolved problems in spectral graph theory,
\emph{Oper. Res. Trans.} \textbf{27} (2023), no.~4, 33--60.

\bibitem{MS}
T.~S.~Motzkin and E.~G.~Straus,
Maxima for graphs and a new proof of a theorem of Tur\'an,
\emph{Canad. J. Math.} \textbf{17} (1965), 533--540.

\bibitem{Nikiforov}
V.~Nikiforov,
Some inequalities for the largest eigenvalue of a graph,
\emph{Combin. Probab. Comput.} \textbf{11} (2002), 179--189.

\bibitem{SMB}
N.~Shaked-Monderer and A.~Berman,
\emph{Copositive and Completely Positive Matrices},
World Scientific, Singapore, 2021.

\bibitem{Wilf}
H.~S.~Wilf,
Spectral bounds for the clique and independence numbers of graphs,
\emph{J. Combin. Theory Ser. B} \textbf{40} (1986), 113--117.

\bibitem{ZZ}
J.~Zeng and X.-D.~Zhang,
A note on the Bollob\'as--Nikiforov conjecture,
\emph{Linear Algebra Appl.} \textbf{710} (2025), 230--242.

\bibitem{Zhang}
S.~Zhang,
On the first two eigenvalues of regular graphs,
\emph{Linear Algebra Appl.} \textbf{686} (2024), 102--110.

\end{thebibliography}
\end{document}
